\documentclass[10pt,twocolumn]{article}

\usepackage[utf8]{inputenc}
\usepackage[T1]{fontenc}
\usepackage{lmodern}
\usepackage[margin=0.75in,top=1in,bottom=1in]{geometry}
\usepackage{amsmath,amssymb}
\usepackage{booktabs}
\usepackage{graphicx}
\usepackage{xcolor}
\usepackage{hyperref}
\usepackage{microtype}
\hypersetup{colorlinks=true,linkcolor=blue!60!black,citecolor=green!50!black,urlcolor=blue!60!black}

\usepackage{../shared/cerebrum-macros}
% Unified notation table for CEREBRUM papers
% Resolve naming conflicts: TSC uses \eta_T for temporal weight; CSA uses \eta for the same concept.
% All papers should import this file and use the macros below for consistency.

\newcommand{\etaT}{\eta_T}        % TSC temporal weight (renamed from generic \eta to avoid conflict)
\newcommand{\etaCSA}{\eta}        % CSA temporal decay weight (10-param formula, position 7)
\newcommand{\alphaCSA}{\alpha}    % CSA semantic similarity weight
\newcommand{\betaCSA}{\beta}      % CSA community score weight
\newcommand{\gammaCSA}{\gamma}    % CSA edge-type weight
\newcommand{\deltaCSA}{\delta}    % CSA distance penalty
\newcommand{\epsilonCSA}{\epsilon}% CSA hop decay
\newcommand{\zetaCSA}{\zeta}      % CSA PageRank prior weight
\newcommand{\iotaCSA}{\iota}      % CSA node recency weight
\newcommand{\muCSA}{\mu}          % CSA synthesis-density penalty
\newcommand{\thetaCSA}{\theta}    % CSA grounding confidence weight

% Graph notation
\newcommand{\KG}{\mathcal{G}}     % Knowledge graph
\newcommand{\Eset}{\mathcal{E}}   % Edge set
\newcommand{\Vset}{\mathcal{V}}   % Vertex/node set
\newcommand{\Cset}{\mathcal{C}}   % Community set
\newcommand{\csascore}{a(u,v,k)}  % CSA attention score

% Traversal notation
\newcommand{\beam}{B}             % Beam width
\newcommand{\hops}{H}             % Number of hops
\newcommand{\mrr}{\text{MRR}}     % Mean Reciprocal Rank
\newcommand{\hitatk}[1]{\text{H@#1}} % Hits@k


\title{\textbf{Experience-Dependent Structural Plasticity in Knowledge Graphs:\\
Bridge Twins and Causal Discovery via STDP}}

% Author block — shared across all CEREBRUM arXiv papers
\author{
  Bryan Alexander Buchorn \\
  Independent Researcher \\
  Las Vegas, NV, USA \\
  \texttt{bryan.buchorn@gmail.com}
}

\date{May 2026}

\begin{document}
\maketitle

% -----------------------------------------------------------------------
\begin{abstract}
Static Knowledge Graphs are a poor model of evolving knowledge: they cannot
incorporate new information from query experience, and they cannot infer causal
structure from temporal event streams. We present two complementary mechanisms for
making KGs structurally plastic. First, the \textbf{Bridge Twin Engine} implements a
graph-theoretic analog to Long-Term Potentiation (LTP): when a reasoning path
repeatedly crosses a community boundary, a relay node is materialized in the
destination community, short-circuiting the path for future queries. On WebQSP, bridge
synthesis raises H@10 from 16.59\% to 20.84\% (+25.6\%). Second, the
\textbf{STDP Causal Engine} adapts Spike-Timing-Dependent Plasticity to infer
directional causal edges from temporal event streams in $O(1)$ amortized time per event
via Lazy Decay. The two mechanisms are biologically parallel: Bridge Twins enact LTP
(repeated co-activation strengthens structural connections) while the STDP engine enacts
Hebbian temporal learning (co-occurrence timing determines causal direction).
Together they form a self-improving KG that reshapes its topology from reasoning history
while simultaneously discovering new causal structure from live data streams.
\end{abstract}

% -----------------------------------------------------------------------
\section{Introduction}
\label{sec:intro}

The assumption that a Knowledge Graph is a fixed artifact is increasingly untenable.
Production KGs receive continuous updates; reasoning systems accumulate experience;
event streams carry temporal causal signals. Yet most KG reasoning systems treat the
graph as a static substrate and discard all experience once a query is answered.

Biological neural circuits exhibit a fundamentally different property: they are plastic.
Long-Term Potentiation (LTP)~\cite{hebb1949,bipoo1998} strengthens synaptic connections
that fire together, encoding experience into structure. Spike-Timing-Dependent
Plasticity (STDP)~\cite{markram1997,sjostrom2010} refines the LTP rule with temporal
precision: the direction of synaptic strengthening depends on the \emph{order} of
pre- and post-synaptic spikes, enabling directional causal inference.

We apply both principles to KG topology. The Bridge Twin Engine enacts LTP at the
community level: frequently traversed inter-community paths become structurally
reinforced via relay nodes. The STDP Causal Engine enacts Hebbian temporal learning
at the edge level: entity co-occurrence timing determines causal edge direction.

\paragraph{Contributions.}
\begin{enumerate}
  \item Bridge Twin materialization with the potentiation rule~(\ref{eq:potentiation})
        and LTD decay, achieving +25.6\% WebQSP H@10.
  \item Proactive bridge synthesis (GraphBridgeEngine) eliminating the cold-start penalty
        at the cost of pre-load time.
  \item The STDP causal weighting rule~(\ref{eq:stdp}) adapted for KG event streams,
        with four-stage significance filtering.
  \item Lazy Decay: an $O(1)$ optimization for continuous causal pair maintenance,
        enabling enterprise-scale event throughput.
  \item A characterization of the interaction effect: Bridge Twin relay paths
        propagate STDP-inferred causal edges, enabling the two mechanisms to compound.
\end{enumerate}

% -----------------------------------------------------------------------
\section{Background}
\label{sec:background}

\paragraph{Hebbian plasticity.}
Hebb's postulate~\cite{hebb1949}: ``neurons that fire together, wire together.''
Formalized by \citet{bipoo1998} as LTP: repeated co-activation of pre- and
post-synaptic neurons increases synaptic weight. The converse — Long-Term Depression
(LTD) — reduces weights for infrequently co-activated pairs.

\paragraph{Spike-Timing-Dependent Plasticity.}
\citet{markram1997} and \citet{bipoo1998} showed that the \emph{order} of spikes
matters: if the pre-synaptic neuron fires before the post-synaptic neuron, the synapse
is potentiated; if the order reverses, it is depressed. This temporal precision enables
directional causal inference from co-occurrence alone.

\paragraph{Graph plasticity in ML.}
Dynamic graph methods~\cite{kazemi2020} model edge evolution over time but do not
materialize structural shortcuts from reasoning experience. Differentiable graph
structure learning~\cite{franceschi2019} optimizes graph topology via gradients but
requires a training objective. The Bridge Twin Engine requires neither: it materializes
shortcuts from query frequency alone.

% -----------------------------------------------------------------------
\section{Bridge Twin Engine: LTP for Knowledge Graphs}
\label{sec:bridge}

\subsection{The Plasticity Problem}

In the \CEREBRUM{} framework~\cite{buchorn2026flagship}, community boundaries represent
semantic discontinuities: an inter-community edge crosses from one attention head to
another, typically at higher cognitive cost than an intra-community edge. When the beam
search repeatedly crosses the same boundary, it is paying this cost on every query.
LTP analogy: this is a synapse that fires repeatedly but has not yet been potentiated.

\subsection{Potentiation Rule}

A relay node $v'_{twin}$ is materialized for source node $v$ in destination community
$C_{dest}$ when:
\begin{equation}
\mathrm{count}(v \to C_{dest}) \geq n_{min}
\quad\text{and}\quad
\cos(\vec{e}_v, \vec{c}_{dest}) \geq \sigma
\label{eq:potentiation}
\end{equation}
where $\vec{c}_{dest}$ is the community centroid and $\sigma = 0.6$ (default).

Upon potentiation:
\begin{enumerate}
  \item Assign $v'_{twin}$ to $C_{dest}$.
  \item Create a \texttt{BRIDGE\_TWIN} edge $(v, v'_{twin})$ with weight
        $w_{bridge} = \cos(\vec{e}_v, \vec{c}_{dest})$.
  \item Mirror all edges of $v$ that terminate in $C_{dest}$ onto $v'_{twin}$.
\end{enumerate}

\subsection{LTD Decay and Pruning}

Relay confidence decays via:
\begin{equation}
c_t = c_0 \cdot \lambda^{\Delta t}
\label{eq:decay}
\end{equation}
where $\lambda \in (0,1)$ is the decay constant and $\Delta t$ is elapsed time since
last traversal. When $c_t < c_{min}$, the relay is pruned during the REM Cycle,
enforcing topological homeostasis.

\subsection{Proactive Synthesis}

The reactive variant of the engine only materializes bridges after they are traversed.
The \textbf{GraphBridgeEngine} extends this with proactive synthesis: at graph load time,
it identifies weakly connected components and pre-materializes bridges across them,
eliminating cold-start latency for early queries.

\subsection{Evaluation}

\begin{table}[h]
\small\centering
\caption{WebQSP impact of bridge synthesis.}
\label{tab:bridge}
\begin{tabular}{lcc}
\toprule
\textbf{Configuration} & \textbf{H@10} & \textbf{Active Bridges} \\
\midrule
No bridge synthesis         & 16.59\% & 0 \\
GraphBridgeEngine (OPT)     & 20.84\% & 8,300 \\
\midrule
Relative gain & \textbf{+25.6\%} & — \\
\bottomrule
\end{tabular}
\end{table}

% -----------------------------------------------------------------------
\section{STDP Causal Engine: Temporal Learning for KG Streams}
\label{sec:stdp}

\subsection{The Temporal Signal Problem}

Entity co-occurrence in a temporal stream carries causal information. If entity $u$
consistently appears before entity $v$ in event logs, $u$ likely causes or precedes $v$.
Standard KG ingestion discards timestamp information, treating events as unordered.

\subsection{The STDP Weighting Rule}

For entity pair $(u, v)$ with spike times $t_{pre}$ and $t_{post}$:
\begin{equation}
\Delta w_{uv} = \begin{cases}
  A_+ \cdot e^{-\Delta t / \tau_+} & \text{if } \Delta t = t_{post} - t_{pre} > 0 \\
  -A_- \cdot e^{-|\Delta t| / \tau_-} & \text{if } \Delta t < 0 \text{ or reverse order}
\end{cases}
\label{eq:stdp}
\end{equation}

where $A_+, A_- > 0$ are potentiation/depression amplitudes and $\tau_+, \tau_-$ are time
constants. Parameters $A_+ = 0.1, \tau_+ = 20$ms, $A_- = 0.12, \tau_- = 20$ms follow
the Bi-Poo model~\cite{bipoo1998}.

When $w_{uv}$ accumulates above a threshold $w_{thresh}$ with sufficient evidence, a
causal edge $u \xrightarrow{\texttt{CAUSES}} v$ is materialized.

\subsection{Four-Stage Significance Filter}

To prevent spurious potentiation from adversarial or rhythmic event floods, four
rejection criteria are applied in order:
\begin{enumerate}
  \item \textbf{Weight threshold}: $w_{uv} \geq w_{thresh}$.
  \item \textbf{Pairing count}: $n_{pairs} \geq n_{min}$ (minimum evidence).
  \item \textbf{Temporal span}: time between first and last pairing $\geq T_{span}$,
        blocking burst floods.
  \item \textbf{Uniformity}: $\chi^2$ test on inter-spike intervals; rejection if
        distribution is too peaked (indicating rhythmic or adversarial artifact).
\end{enumerate}

In adversarial jitter benchmarks (10K events/second, uniform temporal noise), the filter
maintains false-positive rate $\leq$2\%.

\subsection{Lazy Decay}

Na\"{i}ve periodic decay of $N$ causal pairs is $O(N)$. Lazy Decay defers computation
to access time:
\begin{equation}
w'_{uv} = w_{uv} \cdot \lambda^{(T - t_{last})}
\label{eq:lazy}
\end{equation}
where $T$ is the current global step and $t_{last}$ is the pair's last update timestamp.
This reduces maintenance complexity from $O(N)$ per tick to $O(1)$ per event,
enabling constant sub-millisecond latency regardless of accumulated pair count.

\begin{table}[h]
\small\centering
\caption{STDP engine latency vs accumulated pair count.}
\label{tab:stdp_latency}
\begin{tabular}{lc}
\toprule
\textbf{Pair Count} & \textbf{Latency/event (ms)} \\
\midrule
$10^3$   & 0.08 \\
$10^4$   & 0.08 \\
$10^5$   & 0.09 \\
$10^6$   & 0.09 \\
\bottomrule
\end{tabular}
\end{table}

\noindent Constant latency is the key production property: the engine does not degrade
as the causal pair table grows.

% -----------------------------------------------------------------------
\section{Interaction: Plasticity Compounds}
\label{sec:interaction}

Bridge Twin relay nodes propagate STDP-inferred causal edges. When the STDP engine
materializes $u \xrightarrow{\texttt{CAUSES}} v$ and $v$ has a bridge twin $v'$ in a
downstream community, the causal edge is automatically mirrored to $v'$. This means:
\begin{itemize}
  \item Causal chains that cross community boundaries benefit from relay short-circuits.
  \item Bridge materialization accelerates causal discovery by expanding the effective
        neighborhood where STDP events are recorded.
\end{itemize}

In practice, STDP-derived causal edges are tagged with \texttt{source="stdp"} provenance,
allowing the \CSA{} synthesis-density penalty ($-\mu\sigma_d$) to apply appropriate
skepticism without disabling them entirely.

% -----------------------------------------------------------------------
\section{Conclusion}
\label{sec:conclusion}

We presented two biologically-inspired mechanisms for KG structural plasticity. The
Bridge Twin Engine enacts LTP-style reinforcement of frequently traversed paths,
achieving a measurable +25.6\% recall improvement on WebQSP. The STDP Causal Engine
discovers directional causal edges from event stream timing, maintaining constant
per-event latency at production scale via Lazy Decay. Together, they create a
self-improving KG that learns from reasoning experience (Bridge Twins) and from live
data streams (STDP) — without gradient-based training or labeled supervision.

\section*{Acknowledgments}
The author gratefully acknowledges the use of Claude (Anthropic) as a research assistant
for mathematical formalization, code generation, and manuscript preparation. All
conceptual contributions, architectural decisions, experimental design, and intellectual
claims are solely the author's.

\bibliographystyle{plain}
\begin{thebibliography}{99}

\bibitem{hebb1949}
Hebb, D.O. (1949).
\emph{The Organization of Behavior: A Neuropsychological Theory.}
Wiley.

\bibitem{bipoo1998}
Bi, G.Q., \& Poo, M.M. (1998).
Synaptic modifications in cultured hippocampal neurons.
\emph{Journal of Neuroscience}, 18(24).

\bibitem{markram1997}
Markram, H., et al. (1997).
Regulation of synaptic efficacy by coincidence of postsynaptic APs and EPSPs.
\emph{Science}, 275(5297).

\bibitem{sjostrom2010}
Sjöström, J., \& Gerstner, W. (2010).
Spike-timing dependent plasticity.
\emph{Scholarpedia}, 5(2).

\bibitem{pearl2009}
Pearl, J. (2009).
\emph{Causality: Models, Reasoning, and Inference.}
Cambridge University Press.

\bibitem{kazemi2020}
Kazemi, S.M., et al. (2020).
Representation Learning for Dynamic Graphs: A Survey.
\emph{JMLR}, 21.

\bibitem{franceschi2019}
Franceschi, L., et al. (2019).
Learning Discrete Structures for Graph Neural Networks.
\emph{ICML}.

\bibitem{buchorn2026flagship}
Buchorn, B.A. (2026).
CEREBRUM: Training-Free Knowledge Graph Reasoning via Community-Structured Graph Attention.
arXiv preprint.

\bibitem{buchorn2026report}
Buchorn, B.A. (2026).
CEREBRUM v2.51: Complete Technical Specification.
arXiv preprint [CEREBRUM\_REPORT\_ID].

\end{thebibliography}

\end{document}
